\scanpage{207}

or to the pairs $(\alpha,\beta)$ of elements of $\widehat{\Pi}_{0,3}$
satisfying the loop relation (37)
\[
 (97)=(37)\qquad
 \unsure{(\alpha\rho(\alpha)^{-1})}
 =(\beta\sigma_\infty(\beta)^{-1})\ell_1^{\unsure{\nu}}
\]
(which we will then complete by taking $f=\beta$).

\margin{$\alpha,\beta$ \unsure{determined} \unsure{modulo}
$(\alpha,\beta)\longmapsto
(\ell_0^{\unsure{n}}\alpha,\ell_0^{\unsure{n}}\beta)$.}

Up to changing $\alpha$ to $\ell_0^\nu\alpha=\underline{\alpha}_0$,
these were the conventions of \unsure{§ 44|§ 4 h} (where the formula
\unsure{thus} takes the form
\[
 \underbrace{\underline{\alpha}_0\rho(\underline{\alpha}_0)^{-1}}_{\xi}
 =\ell_0^\nu
 \underbrace{(f\sigma_\infty(f)^{-1})}_{\unsure{\varphi=\eta}}
\]
\unsure{).}

Taking (79) into account, the conditions

\deleted{\uline{2° Case}
\unsure{(98)} $g=1$ i.e. $\unsure{\alpha_0}=1$ i.e. $f=\alpha$,
i.e. $\xi=f\rho(f)^{-1}$ i.e.}

(93) (94) also take the form
\[
 (98)\qquad u_{\alpha,\beta,f}
 =\mu(t\sigma_\infty+u)\tau_\infty^{\alpha_2},
\]
and imply that $\underline{u}=u_{\alpha,\beta,f}$
\deleted{\illegible{1}} normalizes the subalgebra of the matrix algebra
$M_2(\widehat{\mathbb Z})$ \uline{generated} by $\sigma_\infty$
(taking into account the relation
\[
 (99)\qquad \tau_\infty(\sigma_\infty)=-\sigma_\infty
 \qquad\text{in }\mathrm{Gl}(2,\widehat{\mathbb Z}).
\]

If we look at the subgroup of $\mathcal M_{0,3}^{!\sim}$
consisting of the $u$ satisfying the equivalent conditions (93) (94) (98)
\[
 (100)\qquad
 \mathcal M_{0,3}^{!\sim}\left(\tfrac12\right)
 =
 \{u\in\mathcal M_{0,3}^{!\sim}\mid [u,\sigma_\infty]=1\}
 \overset{\sim}{\longrightarrow}\mathcal N_{0,3}^{!\sim}
\]
we find that the composite homomorphism

\scanpage{208}
\mspage{403}

(101)
\diagram{
\begin{tikzcd}[ampersand replacement=\&,column sep=large]
\mathcal N_{0,3}^{!\sim}
 \arrow[r,"\sim"]
 \arrow[rr,bend right=22,"k_{1/2}"']
\&
\mathcal M_{0,3}^{!}(\tfrac12)^\sim
 \arrow[r]
\&
\mathcal M_{0,3}^{!\sim}
 \arrow[r,"{u\longmapsto\underline{u}=(u_{\alpha,\beta,f})}"']
\&
\mathrm{Gl}(2,\widehat{\mathbb Z})'
\\
\unsure{\mathcal N_{0,3}^{!\sim}(\tfrac12)}
\end{tikzcd}
}

\margin{\uline{N.B.} $\mathcal M_{0,3}^{!}(\tfrac12)^\sim$
is stable under $x\longmapsto\sigma_\infty x\sigma_\infty^{-1}$
and under $x\longmapsto\tau_\infty x\tau_\infty^{-1}$
(since $\sigma_\infty$ and $\tau_\infty$
\unsure{commute} with $\sigma_\infty$).}

factors through the said normalizer --- if we prefer to view it in terms of the
group $\mathrm{Gl}(2,\widehat{\mathbb Z})\simeq M_2(\widehat{\mathbb Z})^*$,
rather than of the \uline{algebra} of matrices, it factors through
the \uline{normalizer} of the \uline{maximal torus}
\deleted{(\unsure{formed by})}
\[
 (102)\qquad
 \unsure{T_\alpha|T_\sigma}=
 \{t\sigma_\infty+u\quad
 (t,u\in\unsure{\widehat{\mathbb Z}|\mathbb Z},
 \underbrace{t^2+u^2}_{=\det(t\sigma_\infty+u)}
 \in\widehat{\mathbb Z}^{*}).
\]

I suspect that there is \uline{only the} splitting $k_{1/2}$ of (101)
of the extension $\mathcal M_{0,3}^{!\sim}$
(of $\mathcal N_{0,3}^{!\sim}$ by $\widehat{\Pi}_{0,3}$,
\deleted{\unsure{I think}} or of $\mathcal N_{0,3}^{!\sim}$
by $\widehat{\mathcal G}_{0,3}^{+}$)
that factors through the normalizer.

\margin{Surely bogus:
there are also \unsure{all} the curves \uuline{\unsure{isogenous}}
to the elliptic curve with modular invariant $\unsure{j|J}=0$
(\unsure{corresponding} \deleted{in}
the orbit $\{\tfrac12,2,-1\}$ in $U_{0,3}$ under $\unsure{\mathfrak S_3}$).}

\uline{2° Case}\quad This is the case
\[
 (103)\qquad
 g=1\quad\text{i.e. }\unsure{\alpha_0}=1
 \quad\text{i.e. }\alpha=f
 \quad\text{i.e. }\xi=f\rho(f)^{-1},\ \text{i.e.}
\]
\[
 (104)\qquad [u,\rho]=1.
\]

For an arbitrary $u$, there is a
$\gamma\inserted{\in\widetilde{\widehat{\Pi}}_{0,3}
=\widehat{\Pi}_{0,3}\cdot\{1,\tau_\infty\}}$
\deleted{\unsure{unique}} such that
$u'=\mathrm{int}(\gamma)\circ u=\gamma u$
\unsure{commutes|commuting} thus with $\rho$,
namely $\gamma=\unsure{\alpha_0^{-1}}$,
\deleted{\unsure{We thus have}}\inserted{But} we will pay Attention to the fact that
$\gamma$ is not always in $\widehat{\Pi}_{0,3}$, in fact
\[
 (105)\qquad
 \gamma=\unsure{\alpha_0^{-1}}\in\widehat{\Pi}_{0,3}
 \Longleftrightarrow
 \varepsilon_3(\unsure{u})=1
 \quad\text{i.e. }\mu(u)=1\ (3)
\]

\scanpage{209}

If we denote by
$\unsure{\mathcal M_{0,3}^{!\sim}}(\unsure{-\overline{j}})$
the subgroup of $\mathcal M_{0,3}^{!\sim}$
consisting of the $u$ satisfying (103), (104), its image in
$\unsure{\mathcal N_{0,3}^{!\sim}}$ is therefore
the subgroup of index $2$ kernel of
\[
 u\longmapsto
 \underbrace{\varepsilon_3(\mu(u))}_{\unsure{\sim}}
 =\varepsilon_3(u)
\]
i.e.
\[
 \unsure{(106)}\qquad
 \mathcal N_{0,3}^{!\sim}(-\overline{j})
 =\{u\in\mathcal N_{0,3}^{!\sim}\mid
 \varepsilon_3(u)=1\quad
 \text{i.e. }\mu(u)\equiv1\ (3)\}
\]

\margin{\uline{N.B.} $\mathcal M_{0,3}^{!\sim}(-\overline{j})$
is stable under $x\longmapsto\rho x\rho^{-1}$.}

and we have
\[
 \unsure{(107)}\qquad
 \unsure{\mathcal M_{0,3}^{!\sim}}[-\overline{j}]
 \overset{\sim}{\longrightarrow}
 \mathcal N_{0,3}^{!\sim}[-\overline{j}].
\]

Inverting it we thus find a
\deleted{composite homomorphism}\inserted{splitting}

(108)
\diagram{
\begin{tikzcd}[ampersand replacement=\&,column sep=large]
\mathcal N_{0,3}^{!\sim}(\unsure{-\overline{j}})
 \arrow[r,"\sim"]
 \arrow[rr,bend right=20,"{k(-\overline{j})}"']
\&
\mathcal M_{0,3}^{!\sim}(\unsure{-\overline{j}})
 \arrow[r,hook]
\&
\mathcal M_{0,3}^{!\sim}
\end{tikzcd}
}
,
whose composite with $\mathcal M_{0,3}^{!\sim}$ is a homomorphism which
\deleted{\unsure{factors}}\inserted{factors}
through the maximal torus $T_\rho$ of $\mathrm{Gl}(2,\widehat{\mathbb Z})'$,
formed by the
\deleted{$\unsure{c\rho+d}$}
\[
 \unsure{(109)}\qquad
 \{c\rho+d\mid c,d\in\widehat{\mathbb Z},\
 c^2+cd+d^2\in\widehat{\mathbb Z}^{*}).
\]

I suspect that this is the only splitting of
$\mathcal M_{0,3}^{!\sim}$
(or \unsure{even} of $\mathcal M_{0,3}^{!}$)
which has this property. \unsure{This} follows from the equivalent form
of conditions (103) (104), provided by (79):

\margin{\unsure{bogus}, for \unsure{same} reasons as just now\ldots}

\scanpage{210}
\mspage{404}

\[
 (110)\qquad
 \underline{u}=u_{\alpha,\beta,f}=\mu(c\rho+d)
\]
(N.B. $\delta=\varepsilon_3(\unsure{\mu|u})$ is equal to $1$ in this case).

The $u$ in question correspond to the pairs $(\eta,f)$ satisfying
\[
 (111)\qquad \eta\sigma_\infty(\eta)=1
\]
\[
 \unsure{(112)}\qquad
 (f\rho(f)^{-1})=\eta\ell_1^{\unsure{\nu}}
\]

\margin{which we \unsure{will complete|complete} by $\xi=f\rho(f)^{-1}$.}

\margin{We can also write $\beta$ instead of \unsure{$f$},
$f(=\beta)$ \unsure{determined} modulo
$f\longmapsto\ell_0^{\unsure{n}}f$,
\unsure{id.} for $\alpha,\beta$.}

or again to the pairs $(\alpha,\beta)$ satisfying the loop relation
$(97)=(37)$, which we then complete by $f=\alpha$.

\uline{3° Case}\quad It is the one where
\[
 (113)\qquad
 f\equiv1\ (L_0)
 \quad\text{i.e. }\alpha\equiv\alpha_0\ (L_0\text{ on the left})
 \quad\text{i.e. }\beta\equiv\beta_0\ (L_0\text{ on the left})
\]
or again $u(\varepsilon_0)=\varepsilon_0^\mu$
$(\mu\in\widehat{\mathbb Z}^{*})$ i.e.
\[
 (114)\qquad
 [u,\varepsilon_0]=\ell_0^\nu
 \qquad
 (\nu\in\widehat{\mathbb Z}\ -\ \nu=\tfrac{\mu-1}{2}),
\]
\inserted{i.e. $u$ normalizes
$L_0=\{\ell_0^n\mid n\in\widehat{\mathbb Z}\}$ --- or again
(\unsure{enlarging} $\widehat{\Pi}_{0,3}$)
$u$ normalizes $L'_0=\{\varepsilon_0^n\mid n\in\widehat{\mathbb Z}\}$.}

\margin{\uline{N.B.} $L_0=L'_0\cap\widehat{\Pi}_{0,3}$.}

In this case, we \unsure{remove|can remove} the indeterminacy
$(\alpha,\beta,f)\longmapsto
(\ell_0^n\alpha,\ell_0^n\beta,\ell_0^n f)$
by requiring
\[
 (115)\qquad f=1,
\]
which determines $\alpha,\beta$ uniquely.

\deleted{The set}\inserted{group}
$\mathcal M_{0,3}^{!\sim}[0]$ of the $u\in\mathcal M_{0,3}^{!\sim}$
which normalize $L_0$ is in one-to-one correspondence with the set of pairs
$(\alpha,\beta)$
$(\alpha\in\widetilde{\widehat{\Pi}}_{0,3}
=\widehat{\Pi}_{0,3}\cdot\{1,\tau_\infty\},
\ \beta\in\widehat{\Pi}_{0,3})$ satisfying

\margin{We do not \unsure{lift|lift} Therefore \illegible{2}
by $\mathcal M_{0,3}^{!\sim}(0)$,
\illegible{2} $\mathcal M_{0,3}$ and $\mathcal N_{0,3}$
since it \unsure{does not cover|covers only} \illegible{2}
section of the extension \illegible{2} by $\widehat{\Pi}_{0,3}$
--- we \illegible{2} (Proposition \illegible{1})
section $\mathcal M_{0,3}^{!\sim}(0)$ further
\unsure{page 412$^1$}).}

\scanpage{211}

the everlasting loop relation
\[
 \unsure{(115)}\qquad
 (\alpha\rho(\alpha)^{-1})
 =(\beta\sigma_\infty(\beta)^{-1})\ell_1^\nu
 \qquad(\text{for }\nu\in\widehat{\mathbb Z}\text{ suitable})
\]
\inserted{(by $(\alpha,\beta)\longmapsto
\unsure{g|\xi}=\alpha\rho(\alpha)^{-1};
\unsure{h|\eta}=\beta\sigma_\infty(\beta)^{-1}$)}
or again to the pairs $(g,h)$
\inserted{(\unsure{or} $(\xi,\eta)$)}
of elements of $\widehat{\Pi}_{0,3}$ satisfying
\[
 \begin{gathered}
 (116)\quad (=(44)+(46))\\
 g\rho(g)\rho^2(g)=1,\qquad
 \unsure{h}\sigma_\infty(\unsure{h})=1,\qquad
 g=\unsure{h}\ell_1^\nu\quad
 (\nu\in\widehat{\mathbb Z}\text{ suitable}),
 \end{gathered}
\]

\margin{In all this, I drop the ``preliminary
condition'' on the generating relations\ldots}

to the pair \deleted{$(\alpha,\beta)$ \unsure{as}}
$(g,h)$ corresponds the automorphism $u$ defined by
\[
 (117)=(23)\qquad
 u(\sigma_\infty)=\unsure{h}\sigma_\infty,\qquad
 u(\rho)=g\rho
\]
or again by
\[
 (118)\qquad
 \unsure{u}(\sigma_\infty)=\mathrm{int}(\alpha)\sigma_\infty,\qquad
 u(\rho)=\mathrm{int}(\beta)\rho.
\]

In the group $\mathcal M_{0,3}^{!\sim}[0]$, there is the subgroup
\[
 (119)\qquad
 \mathcal M_{0,3}^{!\sim}[0]\cap\widehat{\Pi}_{0,3}\simeq L_0
\]
the inclusion
\[
 (120)\qquad
 \widehat{\mathbb Z}\overset{\ell_0}{\underset{\sim}{\longrightarrow}}
 L_0\lhook\joinrel\longrightarrow\mathcal M_{0,3}^{!\sim}
\]
being described by
\[
 (121)\qquad
 \lambda=\ell_0^n\longmapsto
 (\alpha=\unsure{\lambda|\ell_0^n},
 \beta=\unsure{\lambda|\ell_0^n})
\]
\deleted{\unsure{i.e.}\ \illegible{2}\ $(\alpha=\lambda,\beta=$}

or again by
\[
 (122)\qquad
 \lambda=\ell_0^{\unsure{n}}\longmapsto
 (f=\ell_0^n\ell_1^{-n},g=\ell_0^n\ell_1^{-n})
\]
\inserted{i.e. $f=g=\lambda\sigma_\infty(\unsure{\lambda^{-1}})
=\lambda\rho(\unsure{\lambda})^{-1}$,}
multiplication \deleted{by} on the left by \unsure{the}
$\lambda=\ell_0^n$ in $\mathcal M_{0,3}^{!\sim}[0]$
being \unsure{made explicit|made explicit} by
\[
 (123)\qquad
 (\alpha,\beta)\longmapsto(\ell_0^n\alpha,\ell_0^n\beta)
\]
\[
 (124)\qquad
 (\xi,\eta)\longmapsto
 \ell_0^n\xi\ell_1^{-n},\ \ell_0^n\eta\ell_1^{-n}).
\]

\scanpage{212}
\mspage{405}

Since we have an isomorphism
\[
\mathcal M_{0,3}^{!\sim}[0]/L_0
\xrightarrow{\sim}\mathcal N_{0,3}^{!\sim},
\]
we find the

\uline{Corollary} The group $\mathcal N_{0,3}^{!\sim}$ is in canonical bijection
with the set of pairs
$(\alpha,\beta)\in\Pi_{0,3}^{\sim\wedge}\times\Pi_{0,3}^{\wedge}$
satisfying the loop relation (97)
\[
\alpha\rho(\alpha)^{-1}
=\bigl(\beta\sigma_\infty(\beta)^{-1}\bigr)\ell_1^\nu
\qquad(\text{for }\nu\in\widehat{\mathbb Z}\text{ suitable}),
\]
\uline{modulo} the transformations
$(\alpha,\beta)\mapsto(\ell_0^{\unsure{n|m}}\alpha,\ell_0^{\unsure{n|m}}\beta)$.

The group law \inserted{(of $\mathcal M_{0,3}^{!\sim}(0)$ or of its quotient
$\mathcal N_{0,3}^{!\sim}$)} is made explicit in terms of the $(g,h)$, or of the
$(\alpha,\beta)$, by the ``cocycle formulas'' (52): if $u',u$
\inserted{and $u''=u'\cdot u$} are given by
$\unsure{(g',h')},(g,h),(g'',h'')$, or by
$(\alpha',\beta'),(\alpha,\beta)$ and $(\alpha'',\beta'')$, we will have
\[
\tag{125}
\left\{
\begin{aligned}
g''&=u'(g)g'&h''&=u'(h)h'\\
\alpha''&=u'(\alpha)\alpha'&\beta''&=u'(\beta)\beta'.
\end{aligned}
\right.
\]

\uline{Remark on the notation} $\mathcal M_{0,3}^{!\sim}(q)$,
for $q=\frac12,-\overline j,0$.
In fact for every \uline{algebraic} point
\[
\tag{126}
q\in\mathcal U_{0,3}(\mathbb C)
=\mathbb C\setminus\{0,1\}
=\mathbb P^1_{\unsure{\mathbb C|\mathbb Q}}(\mathbb C)
\setminus\{0,1,\infty\}
\]
we expect \deleted{\unsure{to find}} a partial splitting
\[
\tag{127}
k_q:\unsure{\mathcal N|\mathcal M}_{0,3}^{!\sim}(q)\longrightarrow\mathcal M_{0,3}^{!\sim},
\qquad\operatorname{Im}k_q\overset{\text{definition}}=\mathcal M_{0,3}^{!\sim}(q)
\]
of the extension $\mathcal M_{0,3}^{!\sim}$ of
$\mathcal N_{0,3}^{!\sim}$ by $\Pi_{0,3}^{\wedge}$, or

\margin{In fact, \unsure{in a way} \illegible{1} more
\unsure{precise|precise}, $k_q$ will only be \unsure{determined}
up to an inner \unsure{automorphism}...}

\scanpage{213}

at least \inserted{we will have} a partial splitting over
$\Gamma_{\mathbb Q}(q)$, which we will try to determine ;
$\mathcal N_{0,3}^{!\sim}(q)$ denotes a hypothetical
subgroup of $\mathcal N_{0,3}^{!\sim}$
(which we would have to \unsure{manage} to determine, \unsure{if indeed}
\illegible{1} not too \illegible{1}...), inducing on $\Gamma_{\mathbb Q}$
the \inserted{open} subgroup $\Gamma_{\mathbb Q}(q)$ of
$\Gamma_{\mathbb Q}$ stabilizer of $q$.

\margin{The definition is forced for example when $q$ is
an \unsure{element} of a cyclotomic extension of $\mathbb Q$, in such
a way that $\mathbb Q(q)$ corresponds to a finite discrete
quotient group of $\widehat{\mathbb Z}^{*}$ and $\Gamma_{\mathbb Q}(q)$ is
the kernel of
\diagram{\begin{tikzcd}[ampersand replacement=\&]
\Gamma_{\mathbb Q}\arrow[r,"\chi"]\&\widehat{\mathbb Z}^{*}\arrow[r]\&G
\end{tikzcd}}\ldots}

The analysis of cases $1^\circ$ and $2^\circ$ above in fact amounted to
the determination of such splittings, for $q=\frac12$ and for
$q=-\overline j$ --- in the first case $\Gamma_{\mathbb Q}(q)$ is
the whole of $\Gamma_{\mathbb Q}$, in the second it is the subgroup
of index $2$ kernel of $\varepsilon_2\circ\chi:\Gamma_{\mathbb Q}
\longrightarrow\{\pm1\}$, and idem for
$\unsure{\mathcal N|\mathcal M}_{0,3}^{!\sim}(q)$.
By transforming by $\rho$ \inserted{and $\rho^2$}, we
\deleted{\unsure{would lead}}\inserted{would deduce} \unsure{from the splittings|from the splitting} $k_{1/2}$
the splittings $k_2$ and $k_{-1}$ --- by transforming by any one
of the $\sigma_{\unsure{i|\rho}}$,
\deleted{or by one of the $\unsure{\tau_i|\sigma_i}$}, we would deduce
from the splitting $k_{-\overline j}$ a splitting of type $k_{\unsure{-j|-1}}$.

As for $\mathcal M_{0,3}^{!\sim}[0]$, it corresponds
\inserted{in principle} to the point

\scanpage{214}
\mspage{406}

$0$ of $\widehat{\mathcal U}_{0,3}=\unsure{\mathbb P^1_{\mathbb Q}}$,
there \deleted{\unsure{correspond to it by}}\inserted{two other}
analogous subgroups $\mathcal M_{0,3}^{!\sim}[q]$, with
$q=1$, $q=\infty$, corresponding to the two other points ``at infinity''
of $\mathcal U_{0,3}$, and which are deduced from the preceding one by applying
$\rho$, $\rho^2$.
Here, we have groups which \uline{do} \uline{not} correspond
\inserted{\unsure{strictly}} to \inserted{even} partial sections
of $\mathcal M_{0,3}^{!\sim}$ over $\mathcal N_{0,3}^{!\sim}$,
since the intersection with $\Pi_{0,3}^{\wedge}$ is not reduced to $1$,
but it is one of the three loop groups of type $L_0,L_1,L_\infty$.

But we know that \inserted{the choice of }a \unsure{uniformizer|uniformizer}
of the local ring of $\widehat{\mathcal U}_{0,3}$ at $0$ (say)
determines a splitting of
$\mathcal M_{0,3}^{!\unsure{\sim}}\unsure{(0)|[0]}$, regarded as
an extension of $\mathcal N_{0,3}^{!}$ by $L_0$ --- and later on we will have to
try to grasp \deleted{\unsure{this}}\inserted{such a}
splitting \inserted{(associated with a \unsure{uniformizer})}
from the viewpoint of our calculation, as a \unsure{privileged} way
\deleted{of ``\unsure{normalizing}''}, for \unsure{every}
\deleted{$\alpha,\beta$} \unsure{datum} \inserted{mod $L_0$} of $(\alpha,\beta)$
satisfying a loop relation, of ``normalizing'' the choice of the
$(\alpha,\beta)$ mod $L_0$.
I will come back to this later, surely.

\scanpage{215}
\mspage{407}

\uline{V} \uline{Recapitulation} \uline{in terms of}
$\mathcal M_{0,3}^{!\sim}[0]$

It finally appears that the simplest way to study the loop
automorphisms of $\Pi_{0,3}^{\wedge}$ (commuting outer\textsuperscript{ly}
with $\mathfrak S_3$, i.e. the elements of \deleted{\illegible{1}}
$\mathcal N_{0,3}^{!\unsure{\sim}}$) is to begin with the study of the
subgroup $\mathcal M_{0,3}^{!\unsure{\sim}}[0]$, normalizer
of $L_0$ \inserted{in $\mathcal M_{0,3}^{!\sim}$}.
Setting for $u\in\mathcal M_{0,3}^{!\unsure{\sim}}[0]$
\[
\tag{128}
u(\sigma_\infty)=\eta\sigma_\infty,
\qquad u(\rho)=\xi\rho
\qquad(\eta,\xi\in\Pi_{0,3}^{\wedge})
\]
i.e.
\[
\tag{129}
[u,\sigma_\infty]=\eta\qquad [u,\rho]=\xi
\]
which express the outer commutation of $u$ with $\mathfrak S_3$
(indeed generated by $\sigma_\infty$ and $\rho$), the condition of
normalizing $L_0$ is expressed by a relation of the form
\[
\tag{130}
u(\varepsilon_0)=\varepsilon_0^\mu
\quad\text{or}\quad \unsure{u}(\ell_0)=\ell_0^\mu
\qquad(\mu\in\widehat{\mathbb Z}^{*}).
\]

\margin{(130) is also equivalent, given $[u,\rho]=\xi$, to
\[
\tag{130bis}
\left\{
\begin{aligned}
u(\varepsilon_1)&=\operatorname{int}(\xi)(\varepsilon_1^\mu)\\
&\text{and to}\\
u(\varepsilon_\infty)&=
\operatorname{int}\bigl(\rho^{-1}(\xi^{-1})\,\illegible{1}\bigr)(\varepsilon_\infty^\mu)
\end{aligned}
\right.
\]}

\inserted{Now we have, since $\varepsilon_0=\sigma_\infty\rho$}
\deleted{which, taking $\varepsilon_0=\sigma_\infty\rho$ into account,
takes (by an immediate calculation) the form}
\deleted{\unsure{is written}}
\deleted{(131)} \deleted{$\unsure{\eta\sigma_\infty}$}
\[
u(\varepsilon_0)=u(\sigma_\infty)u(\rho)
=\eta\sigma_\infty\xi\rho
=\eta\sigma_\infty(\xi)\underbrace{\sigma_\infty\rho}_{\varepsilon_0}
=\eta\sigma_\infty(\xi)\varepsilon_0
\]
so (130) can be written
\[
\eta\sigma_\infty(\xi)=\varepsilon_0^{\mu-1}
\qquad\text{i.e., setting }\mu-1=2\nu
\quad(\nu\in\widehat{\mathbb Z})
\]
\[
\eta\sigma_\infty(\xi)=\ell_0^\nu
\qquad\text{i.e., applying }\sigma_\infty
\]
\[
\tag{131}
\sigma_\infty(\eta)\xi=\ell_1^\nu.
\]
Now the relations
\[
\tag{132}
(\eta\sigma_\infty)^2=1,\qquad(\xi\rho)^3=1
\]
can also be written
\[
\tag{133}
\boxed{\eta\sigma_\infty(\eta)=1,\qquad
\xi\rho(\xi)\rho^2(\xi)=1},
\]
whence $\sigma_\infty(\eta)=\eta^{-1}$, so that (131) can
also be written in a more \uline{convenient} way
\[
\tag{134}
\boxed{\xi=\eta\ell_1^\nu}
\qquad(\text{loop relation})
\]

\scanpage{216}

We thus find the

\uline{Theorem} The elements $u$ of
$\unsure{\mathcal M|\mathcal N}_{0,3}^{!\sim}[0]$ --- i.e. the
outer automorphisms of $\Pi_{0,3}^{\wedge}$ which commute outer\textsuperscript{ly}
with $\mathfrak S_3$ and which normalize $L_0$ --- are in one-to-one
correspondence \inserted{(via formulas (128) or 129)}
with the pairs $(\xi,\eta)$ $(\xi,\eta\in\Pi_{0,3}^{\wedge})$,
satisfying relations (133)(134), and such that
\inserted{moreover (``preliminary condition'')}
relations (132) (equivalent : (133)) are
\deleted{\unsure{primitive}}\inserted{generating}
for $\widehat{\mathcal G}_{0,3}^{+}$.
(\deleted{\illegible{1}} In (134) $\nu$ is a suitable element of
$\widehat{\mathbb Z}$, obviously uniquely determined, and subject
to the additional condition $2\nu+1\in\widehat{\mathbb Z}^{*}$)
(which is implied by the preliminary condition, and perhaps implies it...)

\margin{this condition $2\nu+1\in\widehat{\mathbb Z}^{*}$ is in fact
a consequence of equations (133), (134), cf. below, theorem page 409!}

\unsure{And} the obvious inclusion
\[
\tag{135}
\begin{aligned}
L_0&\hookrightarrow\mathcal M_{0,3}^{!\unsure{\sim}}\unsure{(0)|(\omega)}\\
\ell_0^n=\lambda&\longmapsto\operatorname{int}(\lambda)
\end{aligned}
\]
is given by
\[
\tag{136}
\ell_0^n=\lambda\longmapsto
\bigl(\xi=\lambda\rho(\lambda)^{-1}=\ell_0^n\ell_1^{-n},
\quad\eta=\lambda\sigma_\infty(\lambda)^{-1}=\ell_0^n\ell_1^{-n}\bigr)
\]
and more generally the left action of $L_0$ on
$\mathcal M_{0,3}^{!\unsure{\sim}}\unsure{(0)|(\omega)}$ is given,
for $\lambda=\ell_0^n\in L_0$, by
\[
\tag{137}
(\xi,\eta)\longmapsto
\bigl(\xi'=\ell_0^n\xi\ell_1^{-n},\quad
\eta'=\ell_0^n\eta\ell_1^{-n}\bigr)
\]
so
\[
\mathcal M_{0,3}^{!\unsure{\sim}}[0]/L_0
\simeq\mathcal N_{0,3}^{!\unsure{\sim}}
\]
is isomorphic to the quotient of the set of pairs $(\xi,\eta)$
satisfying (133)(134) [and the preliminary condition], by the operations
(137) --- taking into account a \deleted{\unsure{operation}}
natural group structure : to put on this set of pairs
$(\xi,\eta)$, given by

\scanpage{217}
\mspage{408}

\[
\tag{138}
\underbrace{(\xi',\eta')}_{\deleted{\unsure{u'}}}(\xi,\eta)=(\xi'',\eta'')
\]
with
\[
\tag{139}
\xi''=u'(\xi)\xi',\qquad \eta''=u'(\eta)\eta'
\]
$u'$ being the automorphism defined by $(\xi',\eta')$.

\uline{Corollary} Every $u\in\unsure{\mathcal M_{0,3}^{!\sim}}$ can be written \unsure{under|in} the form
\[
\tag{140}
u=f^{-1}u_{\xi,\eta}=\operatorname{int}(f^{-1})\circ u_{\xi,\eta},
\]
with $f\in\widehat{\Pi}_{0,3}$ and $(\xi,\eta)$ satisfying the conditions of the preceding theorem, so that we will have
\[
\tag{141}
u(\unsure{\sigma_\infty})=h\sigma_\infty,\qquad u(\rho)=g\rho,\quad\text{with}
\]
\[
\tag{143}
h=f^{-1}\eta\sigma_\infty(f),\qquad g=f^{-1}\xi\rho(f).
\]
For us to have
\[
f^{-1}u_{\xi,\eta}=f'^{-1}u_{\xi',\eta'}
\]
it \unsure{is necessary and sufficient} that there exist $n\in\mathbb Z$ such that
\[
\tag{144}
\xi'=\ell_0^n\xi\ell_1^{-n},\qquad
\eta'=\ell_0^n\eta\ell_1^{-n},\qquad f'=\ell_0^n f.
\]

\uline{N.B.} \unsure{Thanks} to the loop formula, once the multiplier $\mu$ i.e. $\nu$ is fixed (and in fact $\mu$ is already uniquely determined by the knowledge of $\xi$ alone, as is \unsure{shown} by the adelic calculations in $\mathrm{Gl}(2,\widehat{\mathbb Z})'\ldots$), $\xi$ and $\eta$ determine one another \inserted{by (134)}, so we can eliminate one or the other from the system of equations (133), (134), which becomes either
\[
\tag{145}
\begin{gathered}
\xi\rho(\xi)\rho^2(\xi)=1,\qquad
(\deleted{\unsure{\eta}}\inserted{\xi}\ell_1^{-\nu})\sigma_\infty(\xi\ell_1^{-\nu})=1
\quad\text{i.e.}\\
\sigma_\infty(\xi)=\ell_1^\nu\xi^{-1}\ell_0^\nu
\end{gathered}
\]
or
\[
\tag{146}
\eta\sigma_\infty(\eta)=1,\qquad
(\eta\ell_1^\nu)\rho(\eta\ell_1^\nu)\rho^2(\eta\ell_1^{\unsure{\nu|-\nu}})=1.
\]

\scanpage{218}

The adelic calculation. We again denote by $\xi,\eta$ the images of $\xi,\eta$ in $\mathrm{Sl}(2,\widehat{\mathbb Z})'$, which lift to elements, again denoted $\xi,\eta$, of $\mathrm{Sl}(2,\widehat{\mathbb Z})$.

There is a sign ambiguity, which we remove for $\xi$ by imposing
\[
\tag{147}
\xi\rho(\xi)\rho^2(\xi)=\uuline{+}1
\qquad(\text{in }\mathrm{Sl}(2,\widehat{\mathbb Z}))
\]
\inserted{for $\eta$} there \uline{is no simple way} to remove the \unsure{ambiguity} (if we wish to admit \inserted{provisionally} the possibility of writing in the form $\beta\sigma_\infty(\beta)^{-1}$) \uline{\unsure{of which|thus}} we will have a priori only
\[
\tag{148}
\eta\sigma_\infty(\eta)=\deleted{\unsure{\varepsilon}}\pm1,
\]
with no possibility of changing the sign in the $2^\circ$ case by the change $\eta\mapsto-\eta$ ! Let us set

\margin{\uline{but yes}, we can make the choice of $\eta$ by requiring $\xi=+\eta\ell_1^\nu$ (but this convention \unsure{is} not compatible with a convention \unsure{made} elsewhere $\eta=\beta\sigma_\infty(\beta)^{-1}$).}
\margin{$\in\widehat{\Pi}_{0,3}$}

\[
\tag{148}
\xi=\begin{pmatrix}A&B\\C&D\end{pmatrix},\qquad
\eta=\begin{pmatrix}R&S\\T&U\end{pmatrix}.
\]
We will have
\[
\tag{149}
AD-BC=1\qquad RU-ST=1
\]
and relation (147) is then equivalent to :
\[
\tag{150}
C+\unsure{D}-B=1,
\]
\deleted{\unsure{whereas}} (at least in each $\ell$-adic component we will have to have either this relation, or
\[
\xi_\ell=\begin{pmatrix}A_\ell&B_\ell\\C_\ell&D_\ell\end{pmatrix}
=-\rho^{-1}=\begin{pmatrix}1&-1\\-1&0\end{pmatrix},
\]
\deleted{\unsure{cf.}} proposition on \unsure{page} 394 -- but this solution is incompatible with the condition $\boxed{\xi\equiv1\ (2)}$, whence $A,D\equiv1\ (2)$ and $B,C\equiv0\ (2)\ldots$ so (150) is true in $\mathbb Z_\ell$ for $\forall\ell$, hence in $\widehat{\mathbb Z}\ldots$).

\scanpage{219}
\mspage{409}

Relation (148) can also be written
\[
\eta=\pm\sigma_\infty(\eta^{-1})
\]
now $\sigma_\infty(\eta^{-1})={}^t\eta$, so (148) can be written
\[
\eta=\pm{}^t\eta,
\]
but if we had the sign $-$, we would have $U=-U$ so $2U=0$ so $U=0$, which is incompatible with
\[
\eta\equiv\unsure{\pm1|1}\quad(2)
\]
which implies $R,U\equiv1\ (2)$, $S,T\equiv0\ (2)$. We will therefore have $\eta={}^t\eta$ i.e.
\[
\tag{151}
\eta\sigma_\infty(\eta)=+1,\qquad\text{i.e.}
\]
\[
\tag{152}
S=T.
\]

The loop relation (134) can be written
\[
\tag{153}
\begin{aligned}
\begin{pmatrix}A&B\\C&D\end{pmatrix}
&=\deleted{\begin{pmatrix}1&0\\2\nu&1\end{pmatrix}}
\inserted{\varepsilon}\begin{pmatrix}R&S\\T&U\end{pmatrix}
\begin{pmatrix}1&0\\2\nu&1\end{pmatrix}\\
&=\varepsilon\begin{pmatrix}R+2\nu S&S\\T+2\nu U&U\end{pmatrix}
\qquad\text{with }\varepsilon\in\{\pm1\}
\end{aligned}
\]

\margin{\unsure{Observation} here $\varepsilon=1$, by choosing the sign of $\eta$ suitably (153) (the sign convention does not coincide with a previous convention centered on the $\alpha,\beta$, but from the present viewpoint it seems more \unsure{friendly}).}

and it is equivalent to the relations
\[
\tag{154}
B=\varepsilon S,\qquad D=\varepsilon U\qquad\text{and}
\]
\deleted{\illegible{1}}
\[
\tag{155}
\mu D=1\qquad\text{where }\mu=2\nu+1
\]
\unsure{equivalent|being equivalent} also ; given (154), to
\[
\tag{155bis}
\mu U=\varepsilon,
\]
and the latter merely expresses \inserted{(taking (152) into account)} the relation
\[
\varepsilon\overbrace{(T+2\nu U)}^{A'}
+\overbrace{\varepsilon U}^{B'}-\overbrace{\varepsilon S}^{C'}=1
\]
\unsure{paraphrasing}
\[
A+B-C=1,
\]
it

\scanpage{220}

it merely expresses (in view of (154) and \unsure{(150)} the relation
\[
C=\overbrace{\varepsilon(T+2\nu U)}^{C'}.
\]
Thus the relations (154) (155) express the fact that the two sides of (153) \deleted{\unsure{all} \unsure{their terms}} \inserted{three coefficients} equal \deleted{\illegible{1}} \inserted{\unsure{pairwise}}, \inserted{$\unsure{B=B'},\ C=C',\ D=D'$}, and the equality for the last
\[
A=\overbrace{\varepsilon(R+2\nu S)}^{A'}
\]
then follows from this thanks to the relations \deleted{(149)}, \deleted{which allow us to express}
\[
AD-BC=1,\qquad A'D'-B'C'=1,
\]
which give
\[
A=\frac{1+BC}{D},\qquad A'=\frac{1+B'C'}{D'}
\]
(N.B. we know that $D=D'$ is invertible, by (155)) hence $A=A'$.

We thus find the

\uline{Theorem} For a solution $(\xi,\eta)$ in $\widehat{\Pi}_{0,3}\times\widehat{\Pi}_{0,3}$ of the system of equations (133) (134), again denoting by $\xi,\eta$ elements of $\mathrm{Sl}(2,\widehat{\mathbb Z})$ corresponding to them, the sign ambiguity for $\xi$ being removed by relation (147), we will have the relations (154) and (155), which imply the following consequences:

a) $\mu\inserted{=2\nu+1}$ is invertible, and it is entirely \unsure{determined} by $\xi$ -- it is determined by $\eta$ only up to sign.

\scanpage{221}
\mspage{410}

\uline{Remark}

In view of the relations (149) (150)
\[
AD-BC=1,\qquad \unsure{A|C}+D-B=1,
\]
and of the fact that $D$ is invertible, the matrices $\xi$ are entirely determined by the coefficients $B,D$ --- the pairs $(\xi,\eta)$ of matrices $\in\mathrm{Sl}(2,\widehat{\mathbb Z})$, satisfying
\[
\tag{156}
\left\{\begin{gathered}
\xi\rho(\xi)\rho^2(\xi)=1,\quad
\eta\sigma_\infty(\eta)=1,\quad
\xi=\eta\ell_1^\nu
\qquad(\nu\in\widehat{\mathbb Z}\ \text{\unsure{suitable}})\\
\xi\equiv1\quad(2),\qquad\eta\equiv1\quad(2)
\end{gathered}\right.
\]

\margin{Of course, the introduction of the congruences here is \unsure{entirely} incidental --- the algebraic statement \unsure{is independent} of such conditions...}

are in one-to-one correspondence with the pairs $(B,D)$, $B\in\widehat{\mathbb Z}$, $D\in\widehat{\mathbb Z}^{*}$, \inserted{$B\equiv0\ (2)$}, from which we deduce
\[
\xi=\begin{pmatrix}A&B\\C&D\end{pmatrix},\qquad
\eta=\begin{pmatrix}R&S\\T&U\end{pmatrix},\qquad
\text{and }\nu\in\widehat{\mathbb Z}\text{ by}
\]
\[
\tag{157}
\begin{aligned}
C&=B-D+1,\qquad
A=\frac{1+BC}{D}=\frac1D(1+B^2-BD+B)\\
&&\inserted{{}=\mu(1+B+B^2)-B}
\end{aligned}
\]
\[
\tag{158}
\mu=\frac1D,\qquad\nu=\frac{\mu-1}{2},
\]
\deleted{\unsure{(159)}}
\[
\eta=\xi\ell_1^{-\nu}
=\begin{pmatrix}A&B\\C&D\end{pmatrix}
\begin{pmatrix}1&0\\-2\nu&1\end{pmatrix}
=\begin{pmatrix}A-2\nu B&B\\C-2\nu D&D\end{pmatrix}
\qquad\text{i.e.}
\]

\margin{we have normalized the sign of $\eta$ as \unsure{said} in (156)}

\[
\tag{159}
\text{i.e.}\quad
\unsure{U}=D,\quad S=\unsure{T}=B,\quad
R=A-2\nu B=\mu(1+B^2)=\frac{1+B^2}{D}
\]

Put (compare page~\unsure{398'} (78))
\[
\tag{160}
\underline u_{\xi\deleted{\unsure{,\eta}}}
=\mu\begin{pmatrix}1&B\\0&D\end{pmatrix}
=\begin{pmatrix}\mu&\mu B\\0&1\end{pmatrix}
=\begin{pmatrix}\dfrac1D&B/D\\0&1\end{pmatrix}
\in\mathrm{Gl}(2,\widehat{\mathbb Z})
\]

\margin{\deleted{for the moment we \unsure{have} an element $\xi$ of $\unsure{\widetilde{\widehat{\Pi}}_{0,3}}$, \unsure{which is determined|$\xi$ is determined by} as an element of $\mathrm{Sl}(2,\widehat{\mathbb Z})'$ \unsure{or in|or of} $\mathrm{Sl}(2,\widehat{\mathbb Z})$ \uuline{up to sign}}}

so that knowledge of $\underline u_\xi$ is equivalent to that \inserted{of $B$ and $D$, i.e. to that} of $\xi\in\mathrm{Sl}(2,\widehat{\mathbb Z})$, or again to that of the pairs $(\xi,\eta)$ ($\eta$ normalized by $\eta=+\xi\ell_1^{-\nu}$). (In loc. cit. this element was denoted $u_{\alpha,\beta}$... but it depends only on $\xi\in\widehat{\Pi}_{0,3}$, and even only on $\underline\xi\in\mathrm{Sl}(2,\widehat{\mathbb Z})$ --- and knowledge of it is equivalent to that of $\underline\xi\in\mathrm{Sl}(2,\widehat{\mathbb Z})$ ---)

\margin{\uline{N.B.} We deduce $\xi,\eta$ from knowledge of $\underline u$ by the formulas
\deleted{\illegible{1}}
\[
\tag{161}
\left\{\begin{aligned}
\underline u\cdot{}^t\underline u&=\mu\eta\\
\mu&=\det\underline u
\end{aligned}\right.
\]
whence
\[
\tag{162}
\left\{\begin{aligned}
\eta&=\mu^{-1}\underline u\cdot{}^t\underline u=\mu^{-1}\unsure{,}\\
\xi&=\eta\ell_1^{-\nu}=\mu^{-1}(\underline u\cdot{}^t\underline u)\ell_1^{-\nu}
\end{aligned}\right.
\]
}

\scanpage{222}

\deleted{\unsure{the map}} We will have
\deleted{$\unsure{\mathcal M_{0,3}^{!\sim}(0)}$}
\deleted{\unsure{i.e. $\eta=(\underline u,\sigma_\infty)$}}
\[
(163)\qquad
\underline u_\xi(\sigma_\infty)
=\eta\sigma_\infty
=(\xi \ell_1^{-\nu})\sigma_\infty,
\qquad
\underline u_\xi(\rho)=\xi\rho
\]
\[
(164)\qquad
\underline u_\xi(\varepsilon_0)=\varepsilon_0^\mu
\qquad\text{and}\qquad
u_\xi(\ell_0]=\ell_0^\mu
\]

\margin{The formulas (163) can also be written
\[
\unsure{(\underline u,\sigma_\infty)|[\underline u,\sigma_\infty]}=\eta,
\quad [\underline u,\rho]=\xi
\qquad (163\ \mathrm{bis}),
\]
they give the way to \unsure{reconstruct}
$(\unsure{\xi},\eta)$ using $\underline u_\xi=\underline u$.}

If we consider an arbitrary invertible matrix
\[
\underline u=\begin{pmatrix}L&M\\N&P\end{pmatrix}
\in\mathrm{Gl}(2,\widehat{\mathbb Z})
\]
the relation
\[
\deleted{(\illegible{1})}\qquad
\underline u(\varepsilon_0)=\varepsilon_0^\mu
\quad\text{i.e.}\quad
u\varepsilon_0=\varepsilon_0^\mu
\qquad(\text{where }\mu\in\widehat{\mathbb Z}^{*})
\]
\deleted{\unsure{is equivalent to}}\inserted{is written}
\[
\begin{pmatrix}L&-L+M\\N&-N+P\end{pmatrix}
=\begin{pmatrix}L-\mu N&M-\mu P\\N&P\end{pmatrix}
\]
and is equivalent to: \deleted{\unsure{$N$}}
\[
N=0,\qquad L=\mu P
\]
i.e.
\[
\underline u=\begin{pmatrix}\mu P&M\\0&P\end{pmatrix},
\qquad\text{whence}\quad\det u=\mu P^2
\]

the additional condition
\[
\unsure{(165)}\qquad \det\underline u=\mu
\]
gives the condition
\[
\deleted{(\illegible{1})}\qquad P^2=\underline 1,
\]
the most \unsure{obvious} solution being $P=1$,
which gives \unsure{just} the matrices
\[
\underline u=\begin{pmatrix}\mu&M\\0&1\end{pmatrix}
=\underline u_\xi
\]

\scanpage{223}
\mspage{411}

where $\xi=\begin{pmatrix}A&B\\C&D\end{pmatrix}$
\deleted{\unsure{satisfies}} is \unsure{given|given} in terms of
$B,D$ by $D=\dfrac1\mu$, $B=DM=M/\mu$.

\uline{Theorem} The map
\[
(165)\qquad
\begin{array}{rcl}
\mathcal M_{0,3}^{!\sim}[0]&\longrightarrow&
\mathrm{Gl}(2,\widehat{\mathbb Z})^{\deleted{\unsure{\prime|*}}}\\
\underset{u}{u_{\xi,\eta}}&\longmapsto&\underline u_\xi
\end{array},
\quad\text{or again }u\longmapsto\underline u
\]
$\underline u_\xi$ is given by (160)) is a group homomorphism.

\margin{\uline{N.B.} Since we \unsure{have restricted ourselves to}
$\mathcal M_{0,3}^{!\sim}(0)$, we get a homomorphism into
$\mathrm{Gl}(2,\widehat{\mathbb Z})$,
not just into $\mathrm{Gl}(2,\widehat{\mathbb Z})/\mathord\pm1$.}

\margin{\uline{N.B.} This map factors through the
triangular \unsure{subgroup|subgroup} $\mathcal B$ below (169)}

\uuline{Proof}
Let us consider $u=u_{\xi,\eta}$,
\deleted{$u'=u_{\xi',\eta'}$, we will have \unsure{let}
$u''=u_{\xi'',\eta''}$,}
$\underline u$ is its image in $\mathrm{Gl}(2,\widehat{\mathbb Z})$,
and if $x\in\widehat\Pi_{0,3}$, $\underline x$
\deleted{$\xi''=u'(\xi)\xi'\qquad(\text{in }\widehat\Pi_{0,3})$}
its image in $\mathrm{Sl}(2,\widehat{\mathbb Z})$:
\deleted{we will have \unsure{therefore in}}
\[
(166)\qquad
\underline{u(x)}=\underline u(\underline x)
\left(\stackrel{\text{definition}}{=}
\underline u\,\underline x\,\underline u^{-1}\right)
\]
for it suffices to see it for $x=\sigma_\infty$ and for $x=\rho$,
where this formula reduces \unsure{to} (161). This noted,
if $u=u_{\xi,\eta}\inserted{=u_\xi}$,
$u'=u_{\xi',\eta'}\inserted{=u_{\xi'}}$
are two elements of $\mathcal M_{0,3}^{!\sim}(0)$,
$u''=u_{\xi'',\eta''}=u_{\xi''}$ their product, we will have
\inserted{(139)}
\[
\deleted{\unsure{(167)}}\qquad\xi''=u'(\xi)\xi'
\]
whence
\[
\underline\xi''=\underline{u'(\xi)}\,\underline\xi^{\unsure{\prime}}
=\underline u'(\underline\xi)\,\underline\xi',\quad\text{i.e.}
\]
\[
\deleted{\unsure{(167)}}\quad(167)\qquad
\underline\xi''=\operatorname{int}
(\underline u_{\underline\xi'})(\underline\xi)\mathbin{.}\underline\xi'.
\]
and it suffices to verify that if three solutions
\deleted{\illegible{1}}
\deleted{$\xi,\xi',\xi''$ of the equations}
\[
\xi=\begin{pmatrix}A&B\\C&D\end{pmatrix},
\quad\xi'=\begin{pmatrix}A'&B'\\C'&D'\end{pmatrix},
\quad\xi''=\begin{pmatrix}A''&B''\\C''&D''\end{pmatrix}
\]
\unsure{of the} equations
$AD-BC=1$, $C+D\unsure{-}B=1$, $D\in\widehat{\mathbb Z}^{*}$,
are related by (167), then we have
\[
(168)\qquad
\underline u_{\xi''}=\underline u_{\xi'}\mathbin{.}\underline u_\xi
\]

\scanpage{224}

\inserted{Since $\xi$ and $\underline u=u_\xi$ determine each other
(163bis) it}
\deleted{\illegible{1}} amounts to the same thing to say that if,
on the group
\[
(169)\qquad
\begin{array}{ll}
\mathcal B\subset\mathrm{Gl}(2,\widehat{\mathbb Z})&\\
\Vert&\\[-2pt]
\left\{\begin{pmatrix}\mu&M\\0&1\end{pmatrix}
\ \middle|\ \mu,M\in\widehat{\mathbb Z},\ \mu\text{ invertible}\right\}&
\end{array}
\]
we define (cf. (162))
\[
(170)\qquad
\left\{\begin{aligned}
F_\xi &: \mathcal B\longrightarrow\mathrm{Sl}(2,\widehat{\mathbb Z})\\
F_\xi(\underline u)
&=(\underbrace{\det\underline u}_{\mu})^{-1}
\underbrace{\deleted{\unsure{\underline u\,{}^t\underline u}}}
_{\inserted{[\underline u,\rho]}}
\overset{\deleted{\unsure{\mu-1=2\nu}}}
{\deleted{\unsure{\varepsilon_1^{-(\det\underline u-1)}|\ell_1^{-(\det\underline u-1)}}}}
\end{aligned}\right.
\]
then we have the 1-cocycle relation
\[
(171)\qquad
F_\xi(\underline u'\underline u)
=\operatorname{int}(\underline u')\bigl(F_\xi(\underline u)\bigr)
\mathbin{.}F_\xi(\underline u').
\]
By multiplicativity of $\det$, we are \unsure{reduced}
\deleted{$\unsure{(\underline u,\underline u'\in\mathcal B)}$}
to seeing that
\[
\underline u\longmapsto[\underline u,\rho]
=\underline u\,\rho(\underline u)^{-1}
\]
satisfies the cocycle relation, which is trivial.

\uline{Corollary} The map
\[
(172)\qquad
\begin{array}{ccl}
\mathcal M_{0,3}^{\unsure{!\sim}}
=\mathcal M_{0,3}^{!\sim}[0]\mathbin{.}
\widehat{\mathcal G}_{0,3}^{+}
&\longrightarrow&\mathrm{Gl}(2,\widehat{\mathbb Z})'\\
\deleted{\text{product }\tfrac12\text{ direct}}&&\\
u_\xi\mathbin{.}a&\longrightarrow&(\underline u_\xi)'\mathbin{.}\underline a
\end{array}
\]
where $(\underline u_\xi)'$ \unsure{is the} reduction of $\underline u_\xi$
modulo $\pm1$, and where $\underline a$ denotes the image of $a$ in
$\mathrm{Sl}(2,\widehat{\mathbb Z})'$) is a group homomorphism.

Follows at once from the Theorem and the formulas (163).

\margin{careful, this is \uuline{not} a semidirect product,
since $\mathcal M_{0,3}^{!\sim}(0)\cap
\widehat{\mathcal G}_{0,3}^{+}
=L_0\ne\{1\}$, we must therefore show that the defining formula
does not depend on the way of writing a $u$ in the form
\unsure{$u_\xi a$} --- this amounts to seeing that for
$\lambda\in L_0$, \unsure{these two} images \illegible{2} \unsure{as an element of one or the other}
\unsure{factor}, are the same.}

\scanpage{225}
\mspage{412}

From now on I am going to \unsure{specify|specify} the pair
$(\underline\xi,\underline\eta)$ from the adelic point of view
$(\text{in }\mathrm{Gl}(2,\widehat{\mathbb Z}))$ using
$\underline u_\xi$, which is therefore of the form
\[
(173)\qquad
\underline u=\begin{pmatrix}\mu&M\\0&1\end{pmatrix}
\]
\inserted{with $M\equiv0\ (2)$ if $\underline u$ comes from
$\mathcal M_{0,3}^{!\sim}(0)$}
from which we get $\xi,\eta$ by the formulas \unsure{(163 bis)|(167 bis)},
\[
(174)\qquad
\underline\xi=[\underline u,\underline\rho],\qquad
\underline\eta=[\underline u,\underline\sigma_\infty].
\]

For $\underline u\in\mathcal M_{0,3}^{!\sim}(0)$, defined by
$(\xi,\eta)$, as long as we do not bring in the
possibility of writing $\xi$ and $\eta$ uniquely as
\[
\xi=\alpha\rho(\alpha)^{-1},\quad
\eta=\beta\sigma_\infty(\beta)^{-1}
\]
with $\alpha\in\unsure{\widetilde{\widehat\Pi}_{0,3}|\widehat{\widetilde\Pi}_{0,3}}$,
$\beta\in\widehat\Pi_{0,3}$, the only adelic invariants that we
can associate to it directly in a natural way is the pair
$(\underline\xi,\underline\eta)$, or again \unsure{(what} amounts to the
same thing, but is prettier) $\underline u$.
The elements \deleted{\illegible{1}}
\[
\lambda=\ell_0^n\in L_0\hookrightarrow\mathcal M_{0,3}^{!\sim}[0]
\]
correspond to
\[
\begin{aligned}
(175)\qquad
\xi=\eta=\ell_0^n\ell_1^{-n}
&=\begin{pmatrix}\unsure{1}&-2n\\0&\unsure{1}\end{pmatrix}
\begin{pmatrix}1&0\\-2n&1\end{pmatrix}\\
&=\begin{pmatrix}A=\unsure{1+4n^2}&B=-2n\\C=-2n&D=1\end{pmatrix}
\end{aligned}
\]
\margin{\unsure{please} do the \unsure{forgotten} check
which \unsure{is} referred to in the note \unsure{in the margin}
\unsure{just before}\ldots}
\[
(176)\qquad
\underline\lambda=\unsure{\underline{\underline{\ell_0}}^{\,n}|\underline{\ell_0}^{\,n}}
=\begin{pmatrix}1&-2n\\0&1\end{pmatrix}
\]
moreover we have
\[
(177)\qquad
\underbrace{\begin{pmatrix}1&-2n\\0&1\end{pmatrix}}_{\underline{\ell_0}^{\,n}}
\begin{pmatrix}\mu&M\\0&1\end{pmatrix}
=\begin{pmatrix}\mu&M-2n\\0&1\end{pmatrix}
\]
which shows that \deleted{\illegible{1}}
\deleted{the quotient}\inserted{the quotients of the}
subgroups $\mathcal B'_1$ of $\mathcal B$ (169) consisting of the
